\section{Direct reduced scanning and an exact early-stop criterion}
\label{sec:pointed}

The existing scanner computes an unreduced complex and obtains the reduced
rank by dividing the final rank by two.  We implement the pointed complex
directly.  The underlying construction is standard reduced Khovanov theory,
expressed in the local cobordism category used by Bar--Natan's
algorithm~\cite{barnatan2005,barnatan2007}.  Our contributions here are an
implementation compatible with the current scanner, a basepoint placement
rule with an exact frontier guarantee, and a decision-only stopping rule.
The last rule addresses a specific limitation identified in the earlier
reports: a recognizer need not always finish computing the full homology
before proving that its rank exceeds one~\cite{repository}.  The improvements
do not furnish a general quasi-polynomial bound.

\subsection{The quotient at a persistent endpoint}

Cut an edge of the knot diagram into two boundary endpoints, denoted by
$p$ and $q$, and distinguish $p$.  After its first appearance in the scan,
$p$ remains a boundary endpoint until the end.  For crossingless matchings
$a,b$ on the current boundary, the usual basis of
$\operatorname{Hom}(a,b)$ is indexed by dot subsets on the circles of
$a\cup\overline b$.  Exactly one of these circles contains $p$.
Write $x_p$ for a dot on that circle, and put
$A=\F[x]/(x^2)$.

\begin{proposition}[Pointed quotient]
\label{prop:pointedquotient}
On the category of matchings with persistent endpoint $p$, the subspaces
$x_p\operatorname{Hom}(a,b)$ form a two-sided ideal.  The quotient is an
additive category, and subsequent gluing away from $p$ induces an additive
functor on it.  Applying the quotient to a chain homotopy equivalence gives
a chain homotopy equivalence.
\end{proposition}

\begin{proof}
The fixed endpoint supplies a distinguished boundary segment of every
cobordism.  A dot near that segment may be moved through a composite by
the dot-sliding relation.  Thus composing on either side with a morphism
divisible by the boundary dot produces another morphism divisible by the
same boundary dot.  Gluing away from the segment preserves this property.
Composition and gluing therefore descend to the quotient.  They preserve
addition and direct sums.  Finally, the equation expressing a chain homotopy,
$fg-\mathrm{id}=dh+hd$, remains an equation after applying an additive
functor.  The same applies to the equation for the other composite.
\end{proof}

An invertible cancellation performed before the marked endpoint appears
can consequently be retained.  The prefix is computed in the ordinary
category; after attaching the crossing that introduces $p$, the quotient
functor is applied to the resulting complex.  The prefix need not be rebuilt.
This delayed introduction is essential to the frontier policy below.

If an overlay has $k$ circles, the pointed Hom space has dimension
$2^{k-1}$.  We assign $p$ a fresh integer label smaller than every original
label.  The existing circle-numbering convention, which orders circles by
their minimum label, then makes its circle index zero.  A pointed monomial
is encoded by deleting that dot bit.  This halves the maximum coefficient
bit-vector length for the same pair of pointed matchings.  It does not imply
that the whole scan takes half as much time: object multiplicities,
cancellation order, topology compilation, and Python bookkeeping also matter.

\subsection{Compiling the pointed surface evaluation}

The quotient can be evaluated without first forming a full unpointed
morphism.  Let a connected genus-zero surface component have boundary-circle
set $B$.  The ordinary dotted relations evaluate its undotted version as
\begin{equation}
  \sum_{j\in B}\prod_{i\in B\setminus\{j\}}x_i,
  \label{eq:undottedcomponent}
\end{equation}
and its once-dotted version as $\prod_{i\in B}x_i$.  Two dots give zero.
A closed sphere requires exactly one dot, and a component of positive genus
vanishes over $\F$ in this theory.

\begin{lemma}[Marked surface rule]
\label{lem:markedsurface}
If $p\in B$, the pointed evaluation of this component is
\[
 \begin{cases}
  \displaystyle\prod_{i\in B\setminus\{p\}}x_i,
      &\text{if the component carries no dot},\\[4pt]
  0,  &\text{if it carries at least one dot}.
 \end{cases}
\]
Unmarked components retain their ordinary evaluation.
\end{lemma}

\begin{proof}
Every term of \eqref{eq:undottedcomponent} except the term indexed by $p$
contains $x_p$ and disappears in the quotient.  The once-dotted product
contains $x_p$ as well.  The remaining vanishing relations are unchanged.
\end{proof}

The implementation compiles this rule into the same component plans used
by the original algebra.  Euler characteristics, genus tests, and the number
of closed circles are computed before taking the quotient; deleting a dot
variable must not alter those topological quantities.  Only the dot masks
and the evaluation of the marked component change.  The first marked stage
requires particular care: its incoming masks still use the ordinary basis,
while its outgoing masks use the compressed pointed basis.  Later stages
use compressed masks on both sides.

For a matching with $m$ arcs, the pointed endomorphism ring is
\[
 \F[x_2,\ldots,x_m]/(x_2^2,\ldots,x_m^2).
\]
Its units are precisely the polynomials with constant coefficient one.
Each is its own inverse, since every positive-degree polynomial squares to
zero in characteristic two.  Thus the existing equal-matching pivot tests,
identity shortcuts, and cancellation formulas continue to apply.

After all crossings have been attached, the only endpoints are $p,q$.
There is one matching, and its pointed endomorphism ring is $\F$.
Scalar cancellation therefore leaves a zero differential, whose object
count is the reduced rank.  Closing the cut identifies its marked arc with
the marked circle carrying $A/(x)$.  This is the reduced Khovanov complex.
The alternative convention using $xA$ is isomorphic through $1\mapsto x$,
up to the quantum shift.  The implementation records the same unnormalized
homological degrees as the original scanner and does not output quantum
gradings.

\subsection{A basepoint policy that preserves the frontier}

A straightforward early cut can make the computation substantially worse.
Once both ends of the original edge have been processed, its two cut
endpoints remain permanently on the frontier.  Intermediate matchings then
have extra boundary data, and previously useful cancellations can be delayed.
The number of removed dot variables alone is therefore an inadequate cost
model.

\begin{theorem}[Last-crossing basepoint]
\label{thm:lastedgefrontier}
Fix an order of $n$ crossings and choose an edge whose second occurrence
belongs to the last crossing.  Cut this edge at its two occurrences.  The
cut and original diagrams have equal frontier cardinalities after every
proper prefix of the order.  Their final frontiers have cardinalities two
and zero respectively.  In particular,
\[
 w_{\mathrm{cut}}\leq\max\{w_{\mathrm{original}},2\}.
\]
\end{theorem}

\begin{proof}
Before the first occurrence of the edge, neither frontier contains it or
its replacements.  From the first occurrence until the last crossing, the
original frontier contains precisely that edge label; the cut frontier
contains its first replacement instead.  No other boundary label changes.
At the last crossing the original edge closes, while the second replacement
appears and the first persists.  This also covers a loop whose two
occurrences both lie in the last crossing.
\end{proof}

Among the edges of the last crossing, the implementation selects one whose
first occurrence is earliest.  This keeps the quotient active for as long
as possible subject to the theorem.  It is a deterministic placement policy,
not a proof of optimal running time or minimum live-object count.  An
explicit edge option remains available for controlled experiments.

The distinction was material in development.  On the supplied 36-crossing
stress closure, an initial first-crossing cut required over two million
compositions, whereas the baseline required fewer than twenty-four thousand.
The last-crossing policy instead reduced the composition count below
seventeen thousand.  Section~\ref{sec:experiments} records the work counts
and the mixed timing results.  In particular, the pointed engine remains
opt-in: reduced state sizes do not overcome its bookkeeping costs on every
small input.

\subsection{Safe termination from isolated summands}

The full-rank API continues until the reduced homology has been computed.
A separate decision API can sometimes stop much earlier.  An object is
\emph{isolated} when its row of the differential and its set of incoming
nonzero differential entries are both empty.  This is stronger than being
a cycle: it identifies a direct summand at the level of complexes.

\begin{theorem}[Isolated-summand lower bound]
\label{thm:isolatedbound}
At any stage of the ordinary-then-pointed scan of a validated knot diagram,
suppose there are $s$ isolated matching objects.  Then
\[
 \dim_{\F}\widetilde{\operatorname{Kh}}(K;\F)\geq s.
\]
This remains valid when the marked endpoint has not yet appeared.
\end{theorem}

\begin{proof}
The partial complex is a direct sum of its $s$ one-term isolated complexes
and the remaining complex.  Subsequent tensoring and gluing with the
unprocessed tangle are additive.  Delooping and cancellation give chain
homotopy equivalences, and the pointed quotient preserves those equivalences
by Proposition~\ref{prop:pointedquotient}.  Thus the completed complex is
chain homotopy equivalent to the direct sum of the completed isolated
complexes and a remaining complex.

Every live object has a concrete realization in the processed part of the
original diagram: transfer extends a partial smoothing, delooping removes
closed circles and duplicates the remaining matching, and Gaussian
cancellation deletes objects without introducing a new matching shape.
This realization lies in the actual processed region, which may have holes;
no common disk for all frontier matchings is required.  Completing this
realized matching with the original unprocessed tangle therefore produces
the reduced Khovanov complex of a nonempty classical link, with grading
shifts.  The marked edge is either
already represented by the persistent endpoint or is still in the
unprocessed complement.  It is consequently present in every completed
summand.  Reduction therefore cannot silently turn this summand into the
empty link.

The reduced Khovanov homology of any nonempty link is nonzero.  One direct
argument uses its graded Euler characteristic: for a link with $\ell\geq1$
components, its reduced Jones specialization at one has absolute value
$2^{\ell-1}$.  This is an \emph{integer} Euler characteristic of dimensions,
including when the coefficient field is $\F$; it is not the residue of that
integer modulo two.  Nonzero Euler characteristic forces nonzero homology.
Each completed isolated complex hence contributes at least one dimension,
and dimensions add under direct sums.
\end{proof}

\begin{corollary}
Two isolated objects certify that the input is knotted.
\end{corollary}

\begin{proof}
The theorem forces reduced rank at least two.  Reduced rank one is exactly
the unknot case, using Khovanov unknot detection and its coefficient-field
consequence as in the existing recognizer~\cite{kronheimermrowka,repository}.
\end{proof}

There is no early unknot conclusion from finding zero or one isolated
object.  Nor is a large unfinished complex a homology lower bound: its
objects can still cancel.  The implementation therefore examines actual
isolation, preserves resource-limit exceptions, and keeps the rank and
decision result types distinct.

The decision engine checks after each intermediate reduction.  It also
inspects the final transfer before performing its scalar cancellation;
isolated objects already present there can avoid that final elimination.
Upon finding two, it returns a knot verdict, lower bound two, and no exact
rank.  The associated record contains the scan order, marked edge, stage,
remaining crossing count, and two object IDs with matchings and homological
degrees.  This is \emph{replay data}, not an independently checkable small
topological certificate.  A verifier must replay the preceding reductions
unless a separate chain-equivalence witness is supplied.

Natural-order sums of two, three, and eight Conway diagrams give explicit
before-closure examples: all stop after the first eleven crossings, leaving
eleven, twenty-two, and seventy-seven crossings respectively.  The existing
visible-factorization step already handles this family efficiently.  These
examples establish that the decision mechanism works without computing full
homology; they do not show a new advantage over the complete prior pipeline
on that family.  On the supplied prime examples, isolation generally becomes
visible only at the final transfer.

\subsection{Verification and limitations}

The algebra tests compare the compressed evaluation with ``lift to the
ordinary algebra, compose, then quotient'' for every polynomial input on all
triples of noncrossing matchings with at most six endpoints.  Diagram tests
compare 2,898 exhaustively generated signed braid closures with an independent
dense reduced-cube oracle.  Further tests vary every basepoint edge and three
crossing orders on the smaller diagrams, verify the frontier theorem,
check the differential square after each stage, and replay an early result.
The separate decision API is checked against all 2,898 independent ranks.
Section~\ref{sec:experiments} gives the full counts and reproducible commands.

Cross-stage caches require an additional invariant.  Unpointed topology
plans can be shared because they do not depend on the quotient, but a
numerical morphism value has different bit meanings before and after
compression.  The implementation clears shared numerical results whenever
the pair ``marked on input, marked on output'' changes.  It also uses a
distinct transformed-plan memo within each stage.  Without this distinction,
a numerically identical integer from an earlier cache could represent a
different polynomial.

The pointed Hom spaces still have exponential dependence on frontier width,
and reduced homology can itself have exponentially many generators.  Neither
the quotient nor the frontier theorem bounds the multiplicities of the
intermediate complexes.  The early-stop theorem removes the obligation to
finish some large-rank computations, but supplies no universal bound on when
isolated summands appear.  Proving such a bound for a substantial class,
or finding larger direct summands with efficiently certified nonzero
closures, is a precise next research problem.
